\documentclass[11pt]{article}
\usepackage[utf8]{inputenc}
\usepackage[T1]{fontenc}
\usepackage{amsmath,amssymb}
\usepackage{graphicx}
\usepackage{booktabs}
\usepackage{hyperref}
\usepackage[margin=1in]{geometry}
\usepackage{natbib}
\usepackage{authblk}

\title{Stable Population Structure in Europe Since the Iron Age, Despite High Mobility}

\author[1]{Author A}
\author[1,2]{Author B}
\author[2]{Author C}
\affil[1]{Department of Archaeogenetics, Max Planck Institute for Evolutionary Anthropology, Leipzig, Germany}
\affil[2]{Department of Human Genetics, University of Vienna, Vienna, Austria}

\date{}

\begin{document}
\maketitle

\begin{abstract}
Ancient DNA studies have revealed dramatic population turnovers during European prehistory, including the Neolithic farmer expansion and the Bronze Age steppe migrations. However, the genetic history of Europe during the historical period (from the Iron Age to the present, approximately the last 3,000 years) remains poorly characterized, particularly for regions such as France and Armenia where no whole-genome ancient DNA data existed. Here, we present whole-genome data from 204 historical-period individuals from across Europe and the Mediterranean, including the first genomes from several regions. We find that at least 6.9\% of individuals (14/204) carry ancestry uncommon in their sampling region, indicating substantial cross-Mediterranean and long-distance contacts. Despite this evidence of high individual mobility, population structure analysis reveals remarkable stability from the Iron Age to the present: the correlation between genetic and geographic distance remains consistently high (Mantel $r = 0.78$--$0.83$) across all time periods, and pairwise $F_\text{ST}$ values between ancient and modern populations within regions are uniformly low (0.002--0.007). These findings reconcile the apparent paradox of high individual mobility with population-level genetic stability, suggesting that migrants, while detectable, were sufficiently rare relative to census population sizes that they did not perturb the overall geographic patterning of European genetic variation established by the end of the Bronze Age.
\end{abstract}

\section{Introduction}

The genetic history of Europe has been shaped by major demographic events, most notably the Neolithic expansion of farming populations from the Near East beginning approximately 8,000 years ago and the westward migration of steppe-derived pastoralist populations during the Bronze Age, approximately 5,000--4,500 years ago \citep{haak2015massive, allentoft2015population}. These prehistoric migrations profoundly reshaped European genetic structure, with $F_\text{ST}$ shifts on the order of 0.03--0.08 observed in affected regions.

In contrast, the historical period---spanning from the Iron Age (approximately 3,000 years before present, YBP) through the Roman, Medieval, and modern eras---has received far less attention from ancient genomics studies. Historical records document extensive movements of peoples during this period, including Roman expansion, the Migration Period, Viking incursions, and long-distance trade networks \citep{haber2019genetic}. Whether these historically documented movements translated into significant shifts in population genetic structure, or whether European populations remained essentially stable after the Bronze Age, is an open question.

Previous ancient DNA studies of the historical period have been geographically limited and have produced conflicting signals. Some studies have documented genetic continuity within specific regions \citep{antonio2019ancient}, while others have detected admixture events associated with specific historical episodes \citep{margaryan2020population}. Critically, whole-genome data from several major regions---including France and Armenia---have been entirely lacking for the historical period.

Here, we present whole-genome data from 204 individuals spanning the Iron Age to the Medieval period across Europe and the Mediterranean, including the first ancient genomes from France and Armenia. We use population structure analysis to track genetic variation over three millennia and find that, despite evidence of high individual mobility, European population structure has been remarkably stable since the Iron Age.

\section{Related Work}

The foundational ancient DNA studies of European prehistory established the framework of successive population turnovers. \citet{haak2015massive} demonstrated massive migration from the Pontic steppe during the early Bronze Age, and \citet{allentoft2015population} showed that this migration reshaped the genetic landscape of much of Europe. \citet{lazaridis2014ancient} established the three-component ancestry model for modern Europeans (Western Hunter-Gatherers, Early European Farmers, and Steppe pastoralists).

For the historical period, \citet{antonio2019ancient} provided a genomic perspective on ancient Rome, documenting genetic diversity consistent with cosmopolitan immigration but also showing continuity with earlier Italian populations. \citet{margaryan2020population} characterized Viking-age Scandinavian genomes and documented gene flow from the British Isles and Southern Europe. \citet{haber2019genetic} studied Levantine populations across historical periods and found relative stability. However, a pan-European synthesis across the full historical period has been lacking.

The relationship between genetic and geographic distance, often assessed using Mantel tests, has been used to characterize population structure in both modern \citep{novembre2008genes} and ancient contexts. The stability of this relationship over time is a direct measure of whether population structure is being maintained or disrupted.

\section{Methods}

\subsection{Sample collection and ancient DNA extraction}
We obtained skeletal remains from 204 individuals from archaeological sites across Europe and the Mediterranean (Table~\ref{tab:samples}). Samples were collected from France ($n=38$), Armenia ($n=32$), Italy ($n=42$), Iberia ($n=28$), the Balkans ($n=25$), the Eastern Mediterranean ($n=22$), and other regions ($n=17$). Temporal coverage ranged from 3,000 to 400 YBP.

\subsection{Whole-genome sequencing}
DNA was extracted from petrous bones and teeth using established ancient DNA protocols. Libraries were prepared with UDG treatment to reduce deamination damage. Shotgun sequencing was performed on Illumina platforms, and reads were aligned to the human reference genome (GRCh37). Authentication included contamination estimation, damage pattern verification, and sex determination.

\subsection{Population structure analysis}
Principal component analysis (PCA) was performed by projecting ancient individuals onto axes computed from a panel of modern Eurasian populations. ADMIXTURE analysis was run at $K=2$ through $K=12$. Ancestry outliers were identified as individuals whose genetic profile was inconsistent with the majority of individuals from the same geographic region and time period.

\subsection{Temporal stability analysis}
Pairwise $F_\text{ST}$ was calculated between temporal bins within each region (Iron Age, Roman, Medieval, and modern). Mantel tests assessed the correlation between genetic distance ($F_\text{ST}$) and geographic distance for each time period separately.

\begin{table}[h]
\centering
\caption{Sample overview by region.}
\label{tab:samples}
\begin{tabular}{lccc}
\toprule
Region & $n$ & Time range (YBP) & First genomes? \\
\midrule
France & 38 & 2800--500 & Yes \\
Armenia & 32 & 3000--800 & Yes \\
Italy & 42 & 2500--400 & No \\
Iberia & 28 & 2200--600 & No \\
Balkans & 25 & 2800--700 & No \\
Eastern Mediterranean & 22 & 2600--500 & Partial \\
Other & 17 & Various & Mixed \\
\textbf{Total} & \textbf{204} & \textbf{3000--400} & -- \\
\bottomrule
\end{tabular}
\end{table}

\section{Results}

\subsection{Ancestry outliers reveal high individual mobility}
Of 204 individuals, 14 (6.9\%) carried ancestry profiles inconsistent with their sampling region (Table~\ref{tab:outliers}). Eight individuals showed cross-Mediterranean signals (e.g., Near Eastern ancestry in Western European contexts), four exhibited steppe-enriched profiles in southern regions, and two showed other long-distance ancestry components. This indicates that individual mobility across the Mediterranean and beyond was common during the historical period.

\begin{table}[h]
\centering
\caption{Ancestry outlier analysis.}
\label{tab:outliers}
\begin{tabular}{lcc}
\toprule
Category & Count & \% of total \\
\midrule
Local ancestry & 190 & 93.1 \\
Non-local ancestry (outliers) & 14 & 6.9 \\
\quad Cross-Mediterranean signal & 8 & 3.9 \\
\quad Steppe-enriched in south & 4 & 2.0 \\
\quad Other long-distance & 2 & 1.0 \\
\bottomrule
\end{tabular}
\end{table}

\subsection{Population structure mirrors geography across all time periods}
Despite the presence of ancestry outliers, Mantel tests revealed a consistently strong correlation between genetic distance and geographic distance across all time periods (Table~\ref{tab:mantel}). The correlation coefficient ranged from $r = 0.78$ in the Iron Age to $r = 0.83$ in the modern period, with no evidence of disruption during any intervening period. This pattern of isolation by distance, once established, has persisted for at least 3,000 years.

\begin{table}[h]
\centering
\caption{Mantel test results: correlation between genetic ($F_\text{ST}$) and geographic distance.}
\label{tab:mantel}
\begin{tabular}{lccc}
\toprule
Time period & Mantel $r$ & $p$-value & $n$ (populations) \\
\midrule
Iron Age (3000--2200 YBP) & 0.78 & $<$0.001 & 12 \\
Roman (2200--1500 YBP) & 0.81 & $<$0.001 & 14 \\
Medieval (1500--500 YBP) & 0.80 & $<$0.001 & 11 \\
Modern & 0.83 & $<$0.001 & 15 \\
\bottomrule
\end{tabular}
\end{table}

\subsection{Low $F_\text{ST}$ between historical and modern populations within regions}
Pairwise $F_\text{ST}$ values between historical time slices and modern populations within the same region were uniformly low (Table~\ref{tab:fst}), ranging from 0.002 to 0.007. Even the largest values (Iron Age versus modern in the Balkans, $F_\text{ST} = 0.007$) are an order of magnitude smaller than the shifts observed during the Neolithic transition ($F_\text{ST} = 0.03$--$0.08$). This confirms that no major population replacement has occurred in any of these regions during the last 3,000 years.

\begin{table}[h]
\centering
\caption{Pairwise $F_\text{ST}$ between historical time slices and modern populations.}
\label{tab:fst}
\begin{tabular}{lccc}
\toprule
Region & Iron Age vs.\ Modern & Roman vs.\ Modern & Medieval vs.\ Modern \\
\midrule
France & 0.005 & 0.004 & 0.002 \\
Italy & 0.006 & 0.004 & 0.003 \\
Iberia & 0.004 & 0.003 & 0.002 \\
Balkans & 0.007 & 0.005 & 0.003 \\
\bottomrule
\end{tabular}
\end{table}

\section{Discussion}

Our results reveal a striking dichotomy between individual mobility and population-level stability during the European historical period. While nearly 7\% of sampled individuals carry non-local ancestry---indicating active cross-Mediterranean and long-distance contacts---the overall genetic structure of European populations has remained remarkably stable since the Iron Age. The correlation between genetic and geographic distance shows no significant fluctuation across three millennia, and within-region $F_\text{ST}$ values between ancient and modern populations are an order of magnitude smaller than those observed during the prehistoric population turnovers.

This finding resolves an apparent paradox: historical records describe extensive population movements, yet genetic data show continuity. The resolution lies in scale. Individual migrants, traders, soldiers, and slaves were numerous enough to be detectable in our sample of 204 individuals, but they were rare relative to the resident census populations. Unless migration rates are high enough and sustained enough to shift allele frequencies at a population level, individual mobility does not necessarily translate into population-level genetic change \citep{novembre2008genes}.

The first genomes from historical-period France and Armenia fill important gaps in our understanding of European genetic history. Both regions show genetic profiles consistent with their geographic positions and modern populations, further supporting long-term stability.

These results contrast sharply with the prehistoric period, when the Neolithic transition and steppe migrations caused $F_\text{ST}$ shifts of 0.03--0.08 \citep{haak2015massive}. The dramatic population turnovers of European prehistory were followed by three millennia of relative genetic stasis, during which geography became and remained the dominant predictor of genetic structure.

Limitations include the geographic unevenness of our sample and the fact that some regions and periods are represented by small numbers of individuals. Additionally, our focus on autosomal variation may miss patterns detectable in uniparental markers.

\section{Conclusion}

We demonstrate that European population genetic structure has been remarkably stable since the Iron Age, approximately 3,000 years ago, despite clear evidence of high individual mobility. The geographic patterning of genetic variation established by the end of the Bronze Age has persisted to the present day, with historical-period migrations leaving detectable traces at the individual level but not perturbing population-level structure. This finding underscores the distinction between individual mobility and population-level genetic change and provides a comprehensive baseline for understanding European genetic history during the historical period.

\bibliographystyle{plainnat}
\bibliography{references}

\end{document}
